% appendix_ccfh.tex — out-of-sample forecast performance on the CCFH design.
% Mirrors the structure of Crespo Cuaresma, Feldkircher and Huber's (2016) Section 4.3
% paragraph by paragraph, in their register. All numbers are generated by
% app/scripts/paper_figs/ccfh_matrix.mjs from the evaluation-repository results; an
% editorial pass (2026-07-24) verified every quantitative claim against the tables.

\section{Forecast performance on the CCFH design}\label{app:ccfhperf}

Two strands of literature meet in this exercise. The first studies Bayesian shrinkage as
the route to forecasting with vector autoregressions: from the Minnesota prior of
\citet{DoanLittermanSims1984} and \citet{Litterman1986} and the weighted combination of
its components in \citet{simszha1998}, through the demonstration by
\citet{BanburaGiannoneReichlin2010} that, with shrinkage tightened as the system grows,
very large systems forecast as well as small ones, to
\citet{Koop2013} and \citet{CarrieroClarkMarcellino2015} on specification and prior
choices in medium and large systems, \citet{Giannone2015} on setting the shrinkage
hyperparameters by marginal likelihood --- the tuning approach our estimator carries ---
and \citet{Clark2011} on the density side. The second asks whether modeling
international linkages pays: the GVAR of \citet{PesaranSchuermannWeiner2004} and
\citet{DeesDiMauroPesaranSmith2007} was first systematically evaluated out of sample by
\citet{PesaranSchuermannSmith2009}, who report gains over univariate benchmarks that are
present but uneven across variables and countries; \citet{GreenwoodNimmoNguyenShin2012}
extend the evaluation to probabilistic forecasts, \citet{HanNg2011} to a regional panel,
and \citet{DovernFeldkircherHuber2015} ask directly whether joint modeling of the world
economy improves multivariate forecasts. \citet{CrespoFeldkircherHuber2016} ---
henceforth CCFH --- join the two strands: Bayesian shrinkage estimated within a GVAR,
evaluated over a forty-quarter hold-out against univariate and non-shrinkage
alternatives, with the conclusion that the combination --- not either ingredient alone
--- delivers the improvement. Their design has become the reference for the class, and
this appendix adopts it; \citet{ChudikPesaran2016} survey the wider GVAR literature, and
\citet{AdolfsonLindeVillani2007} established the corresponding benchmark-suite tradition
for open-economy policy models.

The evaluation conventions are equally settled. Point accuracy is reported as root mean
squared error relative to a naive univariate benchmark, density accuracy as log
predictive scores --- the measure grounded in \citet{Geweke2001} and surveyed by
\citet{GewekeWhiteman2006}, with \citet{GewekeAmisano2010} on the comparison of
predictive distributions --- and differences in predictive accuracy are
assessed in the tradition of \citet{DieboldMariano1995} with the small-sample correction
of \citet{HarveyLeybourneNewbold1997}. Three questions recur across this literature
without resolution. Horizons beyond one year are rarely examined --- most GVAR
evaluations, CCFH included, stop at four quarters --- although the shrinkage priors at
issue are designed around low-frequency behavior. The respective contributions of the
prior and of the international linkages are confounded in every comparison that changes
both at once, as \citet{PesaranSchuermannSmith2009} and
\citet{CrespoFeldkircherHuber2016} do by construction. Finally, the information in
observable forward-looking data --- policy guidance, the yield curve --- enters none of
the standard specifications, although the hold-outs of this literature span the
effective-lower-bound years in which that information contradicted historical means most
sharply. The exercise below addresses all three: it extends the horizon grid to nine
quarters, holds the prior fixed while removing the linkages, and evaluates anchors built
from each origin's own term structure.

The evaluation design matches CCFH's in all but four respects, following their
Section~4.3 step by step. The initial estimation sample ends in 2003:Q4, and the 40 quarters
2004:Q1--2013:Q4 form the hold-out. Every model is re-estimated at each origin over an
expanding window. Point forecasts are evaluated by root mean squared error and density
forecasts by log predictive scores; all measures are benchmarked against a fifth-order
univariate autoregression. Four differences remain. First, where they consider one- and
four-quarter horizons, we evaluate every horizon up to nine and report $h=1$, $5$ and $9$
($h=1$ denotes the nowcast quarter, so $h=9$ corresponds to $t+9$ under origin-relative
dating). Second, where their panel is the 36-country, seven-variable dataset of
\citet{DovernFeldkircherHuber2015}, ours is the public 28-country \texttt{pesaranData}
panel with the six common variables (real output, inflation, short- and long-term
interest rates, the real exchange rate, and equity prices; total credit is not
available), and two lags rather than their five, which keeps the 40 re-estimations of 28
country models tractable. Third, we summarize by the cross-country median where their
tables report unweighted averages. Fourth, the \texttt{BGVAR} package runs at its default
prior hyperparameters --- with one stated exception: the eigenvalue truncation is raised
from 1.05 to 1.5, without which no SSVS draw survives on data in near-unit-root levels
--- whereas their specification tunes hyperparameters by marginal likelihood; their
published configuration is reproduced with their archived code and data in
Appendix~\ref{app:cred}. One further departure is common to both sides: the trade
weights are the package's 1980--2016 averages, where CCFH restrict weights to
1980--2003 to keep them real-time. The weights enter our leg and the package's leg
identically, so the head-to-head is unaffected, but neither leg is strictly real-time in
the weights. The two exercises therefore answer separate questions: this
one compares the estimators under identical conditions; that one asks whether their
published results can be reproduced.

\begin{table}[t]
\centering
\caption{Point-forecast accuracy on the CCFH design: relative RMSE against the AR(5)
benchmark, median across the 28 \texttt{pesaranData} countries (40 origins, 2004--2013,
re-estimated at every origin). A dagger marks cells whose summed log predictive score
exceeds that of the AR(5). ``+anch.''\ applies the long-run targets of
Section~\ref{sec:anchors}: drift targets for trending log-levels, the recursive
as-of-origin mean for inflation, and, for the interest rates, the forward-rate path
implied by each origin's own yield curve under the expectations hypothesis (see text). No
future information enters any target. The SSVS and MN columns run the \texttt{BGVAR}
package at default hyperparameters; RW is the last-observation random walk, point-scored
only.}
\label{tab:ccfhmatrix}
\footnotesize\setlength{\tabcolsep}{3pt}
\begin{center}
% GENERATED by app/scripts/paper_figs/ccfh_matrix.mjs — do not edit by hand.
\begin{tabular}{lcccccccc}
\toprule
 & G-BVAR & +anch. & no-solve & +anch. & cl.\ GVAR & SSVS & MN & RW \\
\midrule
\multicolumn{9}{l}{\emph{One quarter ahead (the nowcast, $h=1$)}} \\
\midrule
Output & 0.87$^{\dagger}$ & 0.89$^{\dagger}$ & 0.90$^{\dagger}$ & 0.91$^{\dagger}$ & 0.96$^{\dagger}$ & 1.85 & 1.90 & 1.10 \\
Inflation & 1.02$^{\dagger}$ & 1.07 & 1.01$^{\dagger}$ & 1.04 & 1.26 & 1.68 & 1.76 & 1.05 \\
Short rate & 0.95$^{\dagger}$ & 0.99$^{\dagger}$ & 0.99$^{\dagger}$ & 0.96$^{\dagger}$ & 1.60 & 6.24 & 10.50 & 0.97 \\
Long rate & 1.09 & 1.12 & 1.04 & 1.07 & 2.02 & 10.25 & 14.87 & 0.98 \\
Exch.\ rate & 0.99 & 0.97$^{\dagger}$ & 0.98$^{\dagger}$ & 0.96$^{\dagger}$ & 1.37 & 1.55 & 1.67 & 0.99 \\
Equity & 0.96$^{\dagger}$ & 1.05 & 0.95$^{\dagger}$ & 0.98$^{\dagger}$ & 2.37 & 1.57 & 1.67 & 1.01 \\
\addlinespace
\multicolumn{9}{l}{\emph{Five quarters ahead ($h=5$)}} \\
\midrule
Output & 0.75$^{\dagger}$ & 0.74$^{\dagger}$ & 0.76$^{\dagger}$ & 0.73$^{\dagger}$ & 1.85 & 1.21 & 1.00 & 0.98 \\
Inflation & 1.00$^{\dagger}$ & 1.18 & 0.98$^{\dagger}$ & 1.12 & 1.70 & 2.38 & 1.59 & 1.13 \\
Short rate & 0.85$^{\dagger}$ & 0.87$^{\dagger}$ & 0.89$^{\dagger}$ & 0.90$^{\dagger}$ & 2.52 & 3.11 & 3.51 & 0.84 \\
Long rate & 0.96$^{\dagger}$ & 1.14 & 0.90$^{\dagger}$ & 1.11 & 4.75 & 4.34 & 5.40 & 0.80 \\
Exch.\ rate & 0.90$^{\dagger}$ & 0.78$^{\dagger}$ & 0.86$^{\dagger}$ & 0.78$^{\dagger}$ & 2.65 & 0.96 & 0.88 & 0.78 \\
Equity & 0.87$^{\dagger}$ & 0.94 & 0.85$^{\dagger}$ & 0.91$^{\dagger}$ & 2.88 & 1.32 & 1.09 & 0.90 \\
\addlinespace
\multicolumn{9}{l}{\emph{Nine quarters ahead ($h=9$)}} \\
\midrule
Output & 0.64$^{\dagger}$ & 0.62$^{\dagger}$ & 0.62$^{\dagger}$ & 0.62$^{\dagger}$ & 1.50 & 1.15 & 0.87 & 0.93 \\
Inflation & 0.96$^{\dagger}$ & 1.12 & 0.95$^{\dagger}$ & 1.13 & 2.30 & 2.85 & 1.47 & 1.14 \\
Short rate & 0.72$^{\dagger}$ & 0.77$^{\dagger}$ & 0.80$^{\dagger}$ & 0.76$^{\dagger}$ & 1.83 & 2.41 & 2.37 & 0.76 \\
Long rate & 0.85$^{\dagger}$ & 1.06 & 0.78$^{\dagger}$ & 1.05 & 2.16 & 3.60 & 4.18 & 0.72 \\
Exch.\ rate & 0.79$^{\dagger}$ & 0.75$^{\dagger}$ & 0.81$^{\dagger}$ & 0.74$^{\dagger}$ & 2.45 & 0.93 & 0.77 & 0.73 \\
Equity & 0.84$^{\dagger}$ & 0.89 & 0.86$^{\dagger}$ & 0.90 & 4.23 & 1.36 & 1.01 & 0.86 \\
\bottomrule
\end{tabular}

\end{center}
\end{table}

\paragraph{Overall results, one quarter ahead.} Their one-quarter findings are threefold:
most specifications improve upon the naive benchmark in point forecasts; neither
country-specific shrinkage nor international linkages dominates when considered
separately; and combining the two improves forecast performance considerably. The first
panel of Table~\ref{tab:ccfhmatrix} confirms the first finding for our estimator: the
G-BVAR improves upon the benchmark for output (\CmGbvarYHOne{}), the short rate
(\CmGbvarRHOne{}) and equity prices, is at approximate parity for inflation and the
exchange rate (\CmGbvarDpHOne{} and 0.99), and underperforms only for the long rate. The
classical GVAR attains \CmToolboxYHOne{} for output but exceeds unity for the remaining
variables, and the package's ratios at default hyperparameters exceed unity everywhere
--- output \CmSsvsYHOne{} under SSVS, short rates \CmMnRHOne{} under the Minnesota prior
--- consistent with their observation that the conjugate prior without tuned shrinkage
performs least well, here in more extreme form. On the second and third findings our
design permits a sharper decomposition. Their comparison sets a country-specific BVAR
against a GVAR under a diffuse prior, so the prior and the linkages change together; ours
holds the prior fixed and removes only the foreign variables. The G-BVAR and
no-global-solve columns of Table~\ref{tab:ccfhmatrix} differ by at most
\CmStarsNostarsGap{} at any variable and horizon: conditional on the prior, the
international linkages contribute little to point accuracy on this panel. On this
evidence, the improvement they find from combining shrinkage with linkages owes primarily
to the shrinkage, which in their comparison arrives jointly with the linkages. The same
division of labor is why the platform justifies the coupled solve by cross-country
consistency rather than by forecast accuracy (Section~\ref{sec:example}).

\paragraph{Density forecasts.} Their density results confirm their point results, with
the largest improvements among the shrinkage GVARs. The same holds here: the sum over
countries of the G-BVAR's log-score differences relative to the AR(5) is
\CmGbvarYLpsHOne{} for output at the nowcast and \CmGbvarYLpsHNine{} at $h=9$ (the
corresponding cells carry daggers in Table~\ref{tab:ccfhmatrix}); the no-global-solve
column performs comparably; the package's densities are inferior throughout at default
hyperparameters. For the interest rates the anchored columns match the unanchored
densities once the rate targets are consistent with the yield curve --- the design
question taken up below --- while their inflation and equity daggers are lost.

\paragraph{Longer horizons.} At four quarters they find the naive benchmark difficult to
beat, forecast gains less pronounced than at one quarter, and the Sims--Zha prior
strongest in point accuracy --- evidence that the advantage of the random-walk prior with
sum-of-coefficients dummies grows with the horizon. Extending the grid confirms and
strengthens this pattern: the G-BVAR's output ratio declines at every horizon on the
grid, from \CmGbvarYHOne{} at the nowcast to \CmGbvarYHFive{} at $h=5$ and
\CmGbvarYHNine{} at $h=9$, and the short-rate ratio reaches \CmGbvarRHNine{}. The random
walk is strong for interest rates at intermediate horizons --- \CmRwRHFive{} against the
AR(5) at $h=5$ --- but the G-BVAR matches it there and improves upon it at $h=9$, while
the best package column attains \CmMnYHNine{} for output. Their evidence ends at the
horizon at which the prior family's advantage emerges; on the extended grid the G-BVAR's
advantage increases through the final quarter.

\paragraph{Long-run anchors.} This variant has no counterpart in their study and
illustrates the importance of target selection. A first specification anchors every
variable to its recursive as-of-origin sample mean, and the short-rate ratio rises to
roughly twice the benchmark (\CmAnchMeanRHOne{} at the nowcast, \CmAnchMeanRHNine{} at
$h=9$): a mean-reversion target estimated on three decades of declining rates implies a
return to four or five percent, whereas forward guidance over the effective-lower-bound
years indicated that rates would remain near zero.

Replacing the rate targets with the path priced by the market --- the forward rates
implied by each origin's own yield curve under the expectations hypothesis --- removes
almost all of the damage. That path is the information an overnight-indexed-swap (OIS)
curve would provide, and the only real-time version of it in this panel. The anchored
short rate attains \CmAnchRHOne{} at the nowcast, \CmAnchRHFive{} at $h=5$ and
\CmAnchRHNine{} at $h=9$ --- marginally above the unanchored \CmGbvarRHOne{},
\CmGbvarRHFive{} and \CmGbvarRHNine{} on the median --- with a country mean below that of
the unanchored model at the long horizon (\CmAnchRMeanHNine{} against
\CmGbvarRMeanHNine{}) and a higher log score (\CmAnchRLpsHNine{} against
\CmGbvarRLpsHNine{}). At an effective-lower-bound origin the curve is low and flat, so
the anchor holds the forecast down, as the guidance implied. The long rate records the
remaining cost --- \CmAnchLrHNine{} against the unanchored \CmGbvarLrHNine{} at $h=9$ ---
the consequence of extrapolating a two-point curve without a term-premium model, and
inflation, still anchored to the recursive mean, pays the same penalty as the mean-target
rates in kind (\CmAnchDpHFive{} against \CmGbvarDpHFive{} at $h=5$).

The conclusion is that anchors improve forecasts when the target embodies
forward-looking information --- official long-run targets in the live platform,
yield-curve-implied paths for interest rates --- and not when it is a mechanical sample
moment. Anchoring the level of a trending variable to its historical mean fails severely
(output ratios near six), and anchoring rates to their historical mean imposes the
previous interest-rate regime on the forecast. The natural extension in the live platform
is to take its rate anchors from observed OIS curves --- a data requirement, not a
modeling one.

\begin{table}[t]
\centering
\caption{Cross-country differences, following the regional analysis of
\citet{CrespoFeldkircherHuber2016}: median $\{$interquartile range$\}$ of country-level
relative RMSE against the AR(5), output, one quarter ahead. Latin America contains only
Chile on this panel.}
\label{tab:ccfhregion}
\footnotesize\setlength{\tabcolsep}{4pt}
\begin{center}
% GENERATED by app/scripts/paper_figs/ccfh_matrix.mjs — do not edit by hand.
\begin{tabular}{lcccccc}
\toprule
 & G-BVAR & no-solve & cl.\ GVAR & SSVS & MN & RW \\
\midrule
Europe (12) & 0.84\,\{0.12\} & 0.88\,\{0.12\} & 0.94\,\{0.10\} & 1.76\,\{1.03\} & 1.67\,\{0.71\} & 1.01\,\{0.16\} \\
Other dev.\ (5) & 0.94\,\{0.09\} & 0.93\,\{0.04\} & 1.08\,\{0.15\} & 1.98\,\{0.41\} & 1.77\,\{0.53\} & 1.24\,\{0.15\} \\
Em.\ Asia (8) & 0.80\,\{0.04\} & 0.80\,\{0.07\} & 1.05\,\{0.37\} & 1.75\,\{0.73\} & 1.97\,\{0.83\} & 1.30\,\{0.59\} \\
Lat.\ Am.\ (1) & 1.00\,\{0.00\} & 1.08\,\{0.00\} & 1.26\,\{0.00\} & 1.86\,\{0.00\} & 2.21\,\{0.00\} & 1.36\,\{0.00\} \\
MEA (2) & 0.89\,\{0.06\} & 0.96\,\{0.02\} & 0.88\,\{0.05\} & 2.02\,\{0.55\} & 2.35\,\{0.10\} & 1.40\,\{0.40\} \\
\bottomrule
\end{tabular}

\end{center}
\end{table}

\paragraph{Cross-country differences.} Their regional analysis finds that point-forecast
accuracy varies more across regions than across variables: Latin America shows the widest
dispersion and the least accurate forecasts; European and other developed economies are
homogeneous, with tight cross-sectional distributions; and the distributions tighten when
the priors incorporate country-specific shrinkage. Table~\ref{tab:ccfhregion} confirms
the tightening on our side: the G-BVAR's output medians are \CmRegEuYMed{} in Europe and
\CmRegAsiaYMed{} in emerging Asia, with interquartile ranges of \CmRegEuYIqr{} and
\CmRegAsiaYIqr{}, while the package's dispersion is between four and twenty times wider
(\CmRegSsvsEuYIqr{} against \CmRegEuYIqr{} for SSVS in Europe, \CmRegSsvsAsiaYIqr{}
against \CmRegAsiaYIqr{} in emerging Asia). Their Latin American finding cannot be
re-examined on this panel, which contains only Chile (\CmRegClY{}, at parity with the
benchmark); it is covered instead by the reproduction with their own data in
Appendix~\ref{app:cred}, an exercise whose hold-out spans the same 2004--2013 crisis
window.

\begin{table}[t]
\centering
\caption{The platform's own panel: relative RMSE against the random walk, aggregated over
the six headline economies (US, euro area, Japan, UK, Canada, China) and the pilot's four
origins (2016, 2019, 2022 and 2024, all Q1), shipped V3 configuration (seasonal
treatment, COVID-as-missing). A dagger marks cells whose CRPS also beats the random
walk's. ``+anch.''\ applies the shipped long-run anchors. The classical-GVAR column is
diagnostic only: its unrestricted VECM forecasts are explosive on this near-unit-root
panel (the stationarity guard logs, but does not repair, the affected windows). At
$h=8$ and $9$ the 2024 origin's outcomes are not yet observed, so those aggregates
cover three origins.}
\label{tab:ourpanel}
\footnotesize\setlength{\tabcolsep}{3.5pt}
\begin{center}
% GENERATED by app/scripts/paper_figs/ourpanel_matrix.mjs — do not edit by hand.
\begin{tabular}{lcccccccc}
\toprule
 & G-BVAR & +anch. & ctry BVAR & +anch. & AR(5) & AR(1) & RW-drift & cl.\ GVAR \\
\midrule
\multicolumn{9}{l}{\emph{One quarter ahead (the nowcast, $h=1$)}} \\
\midrule
GDP growth & 1.34$^{\dagger}$ & 1.39$^{\dagger}$ & 1.32$^{\dagger}$ & 1.32$^{\dagger}$ & 1.33$^{\dagger}$ & 1.10$^{\dagger}$ & 0.93$^{\dagger}$ & $6\cdot 10^{2}$ \\
Inflation & 0.84$^{\dagger}$ & 0.85$^{\dagger}$ & 0.86$^{\dagger}$ & 0.86$^{\dagger}$ & 0.72$^{\dagger}$ & 1.02$^{\dagger}$ & 1.04$^{\dagger}$ & $4\cdot 10^{2}$ \\
Policy rate & 0.97$^{\dagger}$ & 1.30 & 0.74$^{\dagger}$ & 0.74$^{\dagger}$ & 0.87$^{\dagger}$ & 1.03 & 5.69 & $9\cdot 10^{2}$ \\
\addlinespace
\multicolumn{9}{l}{\emph{Five quarters ahead ($h=5$)}} \\
\midrule
GDP growth & 0.99$^{\dagger}$ & 0.99$^{\dagger}$ & 0.99$^{\dagger}$ & 0.99$^{\dagger}$ & 0.98$^{\dagger}$ & 0.99$^{\dagger}$ & 0.98$^{\dagger}$ & 75.13 \\
Inflation & 0.81$^{\dagger}$ & 0.80$^{\dagger}$ & 0.77$^{\dagger}$ & 0.77$^{\dagger}$ & 0.65$^{\dagger}$ & 0.78$^{\dagger}$ & 0.77$^{\dagger}$ & $4\cdot 10^{2}$ \\
Policy rate & 1.11 & 1.07 & 0.87$^{\dagger}$ & 0.85$^{\dagger}$ & 0.86$^{\dagger}$ & 0.98 & 0.92 & $5\cdot 10^{2}$ \\
\addlinespace
\multicolumn{9}{l}{\emph{Nine quarters ahead ($h=9$)}} \\
\midrule
GDP growth & 1.00 & 1.09 & 0.98$^{\dagger}$ & 0.97$^{\dagger}$ & 0.85$^{\dagger}$ & 0.87$^{\dagger}$ & 0.87$^{\dagger}$ & $4\cdot 10^{2}$ \\
Inflation & 1.08$^{\dagger}$ & 1.08$^{\dagger}$ & 0.73$^{\dagger}$ & 0.69$^{\dagger}$ & 0.52$^{\dagger}$ & 0.64$^{\dagger}$ & 0.66$^{\dagger}$ & $5\cdot 10^{2}$ \\
Policy rate & 1.33 & 1.25 & 1.04 & 0.99$^{\dagger}$ & 0.85$^{\dagger}$ & 1.00 & 0.88$^{\dagger}$ & $1\cdot 10^{3}$ \\
\bottomrule
\end{tabular}

\end{center}
\end{table}

\paragraph{The platform's own panel.} The same protocol, transplanted to the panel the
platform ships --- 38 blocs, the V3 estimation configuration, the random walk as the
naive benchmark --- closes the appendix; Table~\ref{tab:ourpanel} reports the four-origin
pilot the evaluation protocol requires before its full expanding-window run, which goes
to the companion paper. Three observations carry. First, inflation is where the
platform's model earns its keep against the naive benchmarks at short and medium
horizons (\OpGbvarIHOne{} at the nowcast, \OpGbvarIHFive{} at $h=5$), and its density
forecasts improve on the random walk almost everywhere at the nowcast. Second, the
no-coupling country BVAR matches or betters the coupled model in nearly every cell ---
the pesaranData ablation's conclusion, reproduced on the platform's own panel: the
coupling purchases cross-country consistency, not point accuracy. Third, at this pilot's
scale the univariate benchmarks remain the strongest long-horizon point forecasters ---
the AR(5) reaches \OpArFiveIHNine{} for inflation and \OpArFiveGHNine{} for growth at
$h=9$ --- an expected outcome for four origins that straddle the pandemic, and the
question the full expanding window is designed to settle rather than this gate. The
shipped anchors are close to neutral for growth and inflation; on the policy rate they
cost the coupled model at the nowcast (\OpAnchRHOne{} against \OpGbvarRHOne{}) and help
at the long horizon (\OpAnchRHNine{} against \OpGbvarRHNine{}), with a small benefit to
the country model's rates at $h=9$ (\OpSbvarAnchRHNine{} against \OpSbvarRHNine{}). Two
accounting notes: the four blocs whose post-re-sourcing series begin too late for the
early origins (Australia, Denmark, Indonesia, New Zealand) are excluded from the coupled
model's estimation at those origins by a data-availability rule --- none is a headline
economy, so the aggregates in Table~\ref{tab:ourpanel} are unaffected --- and the pilot
is reported for completeness, not as evidence of forecast superiority: the platform's
case rests on the conditional-scenario machinery of Section~\ref{sec:example}, with the
CCFH-design results above showing the estimator competitive where designs are
comparable.

\paragraph{Summary.} Their three principal conclusions carry over, the first two with a
refinement. First, where they conclude that international linkages and locally induced
shrinkage jointly improve forecasts, our controlled comparison attributes the accuracy to
the prior family and assigns the linkages their distinct role: internal consistency of
the cross-country scenario, the property on which the platform rests. Second, where they
find that no single shrinkage prior dominates, the package-at-defaults columns show the
complementary result: at the one-quarter horizon every prior in the package loses to the
univariate benchmark on every variable, so the marginal-likelihood tuning they emphasize
is necessary for the package to compete. Our estimator carries the same prior family with
the same tuning approach. Third, their regional heterogeneity is confirmed, including the
tightening of cross-sectional dispersion under country-specific shrinkage. On their
design and their metrics, the estimator behind this platform delivers the most accurate
point forecasts for output, equity prices and, at most horizons, the short rate and the
exchange rate; the long rate remains with the random walk (\CmRwLrHNine{} at $h=9$
against the family's best \CmGbvarLrHNine{}). Its output advantage widens at every
horizon on the grid, to \CmGbvarYHNine{} at $t+9$ --- a result their four-quarter horizon
could not reveal.
